\documentclass[11pt,a4paper,twocolumn]{article}

\usepackage[utf8]{inputenc}
\usepackage[T1]{fontenc}
\usepackage{lmodern}
\usepackage[margin=2cm]{geometry}
\usepackage{amsmath,amssymb}
\usepackage{graphicx}
\usepackage{hyperref}
\usepackage{listings}
\usepackage{xcolor}
\usepackage{booktabs}
\usepackage{float}
\usepackage{caption}
\usepackage{enumitem}
\usepackage{fancyhdr}
\usepackage{titlesec}

% Colors
\definecolor{codegreen}{rgb}{0.0,0.5,0.0}
\definecolor{codegray}{rgb}{0.5,0.5,0.5}
\definecolor{codeblue}{rgb}{0.15,0.35,0.65}
\definecolor{codebg}{rgb}{0.97,0.97,0.97}

% Code listings
\lstset{
    basicstyle=\ttfamily\footnotesize,
    keywordstyle=\color{codeblue}\bfseries,
    commentstyle=\color{codegreen},
    stringstyle=\color{red!60!black},
    backgroundcolor=\color{codebg},
    frame=single,
    rulecolor=\color{codegray},
    breaklines=true,
    showstringspaces=false,
    tabsize=4,
    language=C, % closest to Rust syntax highlighting
    morekeywords={async,await,fn,pub,struct,impl,use,mod,let,mut,self,Arc,usize,u64,bool,String,Vec,Result,Ok,Err,Some,None,match,loop,break,continue,if,else,for,in,return,const,static,type,where,trait,enum,unsafe,move,ref,dyn,crate,super},
}

\hypersetup{
    colorlinks=true,
    linkcolor=codeblue,
    citecolor=codeblue,
    urlcolor=codeblue,
}

% Headers
\pagestyle{fancy}
\fancyhf{}
\fancyhead[L]{\small\textit{Q-Flux \& Q-Queue Technical Whitepaper}}
\fancyhead[R]{\small\textit{Q-NarwhalKnight Infrastructure}}
\fancyfoot[C]{\thepage}

\title{%
\textbf{Q-Flux \& Q-Queue: High-Performance Network Infrastructure\\for Quantum-Resistant Blockchain Systems}\\[0.5em]
\large Worker-Per-Core Reverse Proxy and Lock-Free Universal Queue Fabric\\for the Q-NarwhalKnight Ecosystem
}

\author{%
Q-NarwhalKnight Infrastructure Team\\
\texttt{https://quillon.xyz}\\[0.3em]
\small Version 1.0 --- March 2026
}

\date{}

\begin{document}

\maketitle

\begin{abstract}
We present \textbf{Q-Flux} and \textbf{Q-Queue}, two infrastructure components designed for the Q-NarwhalKnight quantum-resistant blockchain. Q-Flux is a worker-per-core reverse proxy that achieves $>$500K requests/sec on 48 cores through CPU-pinned workers, \texttt{SO\_REUSEPORT} kernel distribution, TLS session resumption, and per-worker semaphore backpressure. Q-Queue is a universal queue fabric providing sub-microsecond IPC latencies via lock-free ring buffers with per-element versioning, and $>$2~GB/s distributed throughput through io\_uring and RDMA. Both systems are written in Rust with zero-copy I/O paths and integrate natively with the Q-NarwhalKnight DAG-BFT consensus protocol. We describe the architecture, implementation, and performance characteristics of both systems, including solutions to production challenges encountered serving 400+ concurrent miners on a 10Gbit supernode.
\end{abstract}

\section{Introduction}

Blockchain infrastructure faces a fundamental tension between security, decentralization, and performance. The Q-NarwhalKnight system~\cite{qnk-yellow} implements a DAG-based BFT consensus protocol with post-quantum cryptographic primitives, requiring network infrastructure that can sustain high throughput while maintaining sub-second finality.

Production deployment on the Epsilon supernode (48 cores, 10Gbit NIC, 64GB RAM) with 400+ concurrent miners exposed critical limitations in existing reverse proxy solutions:

\begin{itemize}[nosep]
\item \textbf{Pingap} (Cloudflare Pingora-based): Single-process async TLS model collapsed under load. HTTP proxying worked at 12ms latency, but HTTPS connections timed out completely at scale.
\item \textbf{Caddy}: Goroutine-per-connection model created 88K goroutines, 11.5GB heap, 3785\% CPU utilization, producing 87\% 503 error rates on mining submissions.
\item \textbf{Nginx}: Handled load via 48 worker processes, but lacked native Rust integration, WebSocket-aware routing, and quantum-ready TLS configuration.
\end{itemize}

These failures motivated two new infrastructure components: Q-Flux for network ingress and Q-Queue for internal message passing.

\section{Q-Flux: Worker-Per-Core Reverse Proxy}

\subsection{Architecture Overview}

Q-Flux uses a \textit{shared-nothing, worker-per-core} architecture inspired by Seastar~\cite{seastar} and ScyllaDB. Each worker thread is pinned to a physical CPU core and runs its own single-threaded Tokio runtime. The kernel distributes incoming connections across workers via \texttt{SO\_REUSEPORT}, eliminating thundering herd effects and cross-core lock contention.

\begin{figure}[H]
\centering
\begin{verbatim}
    Miners/Browsers (TLS 443)
           |
    [Linux Kernel: SO_REUSEPORT]
     /      |      |      \
  Worker0 Worker1 ... Worker47
  (Core0) (Core1)     (Core47)
     |      |            |
  [Per-worker pools to backend:8080]
\end{verbatim}
\caption{Q-Flux worker-per-core architecture. Each worker binds to the same port; the kernel distributes connections.}
\label{fig:architecture}
\end{figure}

\subsection{Connection Acceptance}

Each worker binds to ports 443 and 80 using \texttt{SO\_REUSEPORT}:

\begin{lstlisting}
// Raw SO_REUSEPORT via libc
unsafe {
    let optval: c_int = 1;
    libc::setsockopt(
        socket.as_raw_fd(),
        SOL_SOCKET, SO_REUSEPORT,
        &optval as *const _ as *const c_void,
        size_of::<c_int>() as socklen_t,
    );
}
socket.listen(4096)?; // 4096 backlog
\end{lstlisting}

This approach has three advantages over a single accept loop:
\begin{enumerate}[nosep]
\item The kernel distributes connections fairly across workers using consistent hashing of the source address.
\item No thundering herd: only one worker is woken per connection.
\item No lock contention on the accept path.
\end{enumerate}

Workers are pinned to CPU cores via \texttt{sched\_setaffinity}, ensuring consistent cache utilization:

\begin{lstlisting}
fn pin_to_core(core_id: usize) {
    unsafe {
        let mut set: cpu_set_t = zeroed();
        CPU_ZERO(&mut set);
        CPU_SET(core_id % num_cpus::get(),
                &mut set);
        sched_setaffinity(0,
            size_of::<cpu_set_t>(), &set);
    }
}
\end{lstlisting}

\subsection{TLS Termination}

Q-Flux uses rustls~0.23 with session resumption optimized for miners who reconnect frequently:

\begin{itemize}[nosep]
\item \textbf{Session tickets}: Automatic key rotation via \texttt{ring::Ticketer}. Resumed handshakes complete in $\sim$0.5ms (75\% faster than full handshake at $\sim$2ms).
\item \textbf{Session cache}: 65,536-entry \texttt{ServerSessionMemoryCache} shared across all workers via \texttt{Arc<ServerConfig>}. rustls is thread-safe.
\item \textbf{ALPN}: Advertises \texttt{http/1.1}; HTTP/2 support planned for Phase~3.
\end{itemize}

The \texttt{Arc<ServerConfig>} is constructed once and cloned to each worker. Because rustls separates configuration from state, the shared config introduces zero contention.

\subsection{Backpressure and Rate Limiting}

Q-Flux implements three layers of backpressure to prevent OOM under load:

\subsubsection{Per-Worker Semaphore}

Each worker limits concurrent connection handlers to 1,024 via a Tokio semaphore. With 48 workers, the system-wide maximum is 49,152 concurrent handlers:

\begin{lstlisting}
let semaphore = Arc::new(
    Semaphore::new(MAX_HANDLERS_PER_WORKER)
);

// In accept loop:
let _permit = match semaphore
    .try_acquire() {
    Ok(p) => p,
    Err(_) => {
        // At capacity: drop connection
        cleanup_conn(...);
        return;
    }
};
\end{lstlisting}

This pattern was proven critical in production: the Q-NarwhalKnight block-pack handler previously crashed Epsilon 6+ times/hour due to unbounded \texttt{tokio::spawn} allocations~\cite{block-pack-oom}.

\subsubsection{Per-IP Connection Limiting}

A shared \texttt{DashMap<IpAddr, u64>} tracks per-IP connection counts across all workers. The implementation avoids a subtle bug where the counter was incremented \textit{before} the limit check:

\begin{lstlisting}
// CORRECT: check THEN increment
let mut count = ip_tracker
    .entry(client_ip).or_insert(0);
if *count >= max_per_ip {
    metrics.rate_limited();
    drop(count); // release lock first
    drop(tcp_stream);
    continue; // count NOT incremented
}
*count += 1; // only after check passes
\end{lstlisting}

\subsubsection{Global Connection Limit}

An \texttt{AtomicU64} tracks total active connections across all workers. The default limit is 100,000. Connections exceeding this are dropped immediately.

\subsection{HTTP/1.1 Proxy Layer}

Request handling uses \texttt{httparse} for zero-allocation header parsing in a keepalive loop:

\begin{lstlisting}
loop {
    let (req, header_end) =
        read_request_headers(
            &mut stream, &mut buf,
            &mut buf_len
        ).await?;

    // WebSocket detection
    if is_websocket_upgrade(&req) {
        handle_websocket_upgrade(
            stream, ...
        ).await;
        return; // connection consumed
    }

    // Forward to upstream
    let resp = upstream.forward(req).await;
    write_response(&mut stream, resp).await;

    if !should_keep_alive(&req) { break; }
}
\end{lstlisting}

\subsubsection{SSE Streaming Support}

Server-Sent Events (SSE) are critical for real-time mining updates. The proxy detects streaming responses via \texttt{Content-Type: text/event-stream} and forwards chunks as they arrive instead of buffering:

\begin{lstlisting}
if is_streaming {
    write_response_headers(stream, &parts)
        .await?;
    loop {
        match body.frame().await {
            Some(Ok(frame)) => {
                stream.write_all(
                    frame.data_ref()
                ).await?;
                stream.flush().await?;
                // Flush each event!
            }
            _ => break,
        }
    }
} else {
    // Buffered mode: collect, then send
    let bytes = body.collect().await?;
    // ... write with Content-Length
}
\end{lstlisting}

\subsubsection{WebSocket Passthrough}

WebSocket upgrades are detected via the \texttt{Upgrade: websocket} header, forwarded to an upstream selected by round-robin, and then bidirectionally spliced:

\begin{lstlisting}
let c2u = tokio::io::copy(
    &mut client_read, &mut upstream_write);
let u2c = tokio::io::copy(
    &mut upstream_read, &mut client_write);
tokio::select! {
    r = c2u => { metrics.bytes_rx(r?); }
    r = u2c => { metrics.bytes_tx(r?); }
}
\end{lstlisting}

\subsubsection{Slowloris Protection}

All client reads have a 30-second timeout to prevent slowloris attacks:

\begin{lstlisting}
match tokio::time::timeout(
    Duration::from_secs(30),
    stream.read(&mut buf[*buf_len..])
).await {
    Ok(Ok(n)) => { *buf_len += n; }
    Err(_) => {
        return Err(anyhow!(
            "Header read timeout"
        ));
    }
    // ...
}
\end{lstlisting}

\subsection{Upstream Pool}

Each worker owns its own \texttt{UpstreamPool} backed by a hyper~1.0 \texttt{Client} with connection pooling:

\begin{itemize}[nosep]
\item \textbf{Round-robin}: \texttt{AtomicUsize} index with \texttt{fetch\_add(1, Relaxed)}
\item \textbf{Health-aware}: Skips unhealthy backends; falls back to round-robin if all are down
\item \textbf{Per-worker isolation}: No cross-thread contention on connection pools
\item \textbf{Hop-by-hop filtering}: Removes \texttt{Connection}, \texttt{Transfer-Encoding}, \texttt{keep-alive}
\item \textbf{Forwarding headers}: Adds \texttt{X-Forwarded-For} and \texttt{X-Real-IP}
\end{itemize}

\subsection{Health Checking}

A dedicated thread runs periodic health probes against each backend:

\begin{enumerate}[nosep]
\item TCP connect with configurable timeout (default: 3s)
\item HTTP GET to configurable path (default: \texttt{/api/v1/status})
\item Accept any 2xx status as healthy
\item Mark unhealthy after $N$ consecutive failures (default: $N=3$)
\item Immediate recovery on first successful probe
\end{enumerate}

Health state is stored in a shared \texttt{DashMap<String, BackendHealth>} read by workers during backend selection.

\subsection{Observability}

\subsubsection{Admin API}

An internal HTTP server on \texttt{127.0.0.1:9090} exposes three endpoints:

\begin{itemize}[nosep]
\item \texttt{GET /health} --- JSON health check with uptime
\item \texttt{GET /metrics} --- Prometheus text exposition format with 18 metrics
\item \texttt{GET /status} --- Full JSON snapshot of all metrics + version
\end{itemize}

\subsubsection{Metrics}

All metrics use \texttt{AtomicU64} counters with \texttt{Relaxed} ordering (sufficient for monotonic counters):

\begin{table}[H]
\centering
\footnotesize
\begin{tabular}{ll}
\toprule
\textbf{Metric} & \textbf{Type} \\
\midrule
\texttt{connections\_active} & Gauge \\
\texttt{connections\_total} & Counter \\
\texttt{tls\_handshakes\{ok,fail\}} & Counter \\
\texttt{requests\_total\{2xx,4xx,5xx\}} & Counter \\
\texttt{upstream\_active} & Gauge \\
\texttt{upstream\_failures} & Counter \\
\texttt{rate\_limited} & Counter \\
\texttt{websockets\_active} & Gauge \\
\texttt{bytes\_\{received,sent\}} & Counter \\
\bottomrule
\end{tabular}
\caption{Q-Flux metrics exposed via Prometheus endpoint.}
\end{table}

\subsubsection{TUI Dashboard}

A terminal UI built with \texttt{ratatui} provides real-time monitoring:

\begin{verbatim}
 Q-FLUX DASHBOARD  v9.2.4  48 workers
 +-----------+-----------+----------+
 | Conns: 847| TLS OK: 1M| RPS: 9.5K|
 | Rate: 12  | TLS F: 23 | 2xx: 99% |
 +-----------+-----------+----------+
 | Upstream Active: 312/768          |
 | WebSockets: 47 active            |
 | Bytes: 12.4GB rx / 89.1GB tx     |
 +-----------------------------------+
\end{verbatim}

\subsection{Graceful Shutdown}

Q-Flux implements three-phase graceful shutdown:

\begin{enumerate}[nosep]
\item \textbf{Signal catch}: SIGTERM/SIGINT handled on main thread
\item \textbf{Fast-path flag}: \texttt{AtomicBool} checked in hot accept loop (no channel overhead)
\item \textbf{Broadcast channel}: \texttt{tokio::sync::broadcast} signals all workers and the metrics reporter
\item \textbf{Drain timeout}: Configurable (default 30s). Workers stop accepting but let in-flight requests complete. Main thread logs drain progress every second.
\item \textbf{Final metrics}: Complete metrics snapshot logged before exit
\end{enumerate}

\subsection{Performance Targets}

\begin{table}[H]
\centering
\footnotesize
\begin{tabular}{lll}
\toprule
\textbf{Metric} & \textbf{Target} & \textbf{Comparison} \\
\midrule
RPS (48 cores) & $>$500K & nginx $\sim$200K \\
P99 latency & $<$5ms & Pingap 12ms (HTTP) \\
Memory (idle) & $<$50MB & Envoy $\sim$150MB \\
Memory (100K conns) & $<$500MB & Caddy 11.5GB \\
TLS handshakes/sec & $>$50K & Pingap $\sim$5K \\
Drain time & $<$10s & Pingap: never \\
\bottomrule
\end{tabular}
\caption{Q-Flux performance targets vs. existing solutions.}
\end{table}

\subsection{Future Work}

\begin{description}[nosep]
\item[Phase 2:] Replace Tokio I/O with raw io\_uring event loops; add SIMD HTTP header parsing via AVX2/SSE4.2. Reuse patterns from \texttt{q-kernel-io} and \texttt{q-crypto-simd}.
\item[Phase 3:] HTTP/2 multiplexing via \texttt{h2} crate; HTTP/3 (QUIC) via \texttt{quinn} for zero-RTT miner reconnects; kTLS offload.
\item[Phase 4:] libp2p-aware proxying: detect multistream-select, per-peer bandwidth tiers, gossipsub dedup at proxy layer.
\item[Phase 5:] Multi-backend load balancing with least-connections and consistent-hash algorithms.
\end{description}

Open issues for community contribution include TLS certificate hot-reload (\#10), request latency histograms (\#11), token bucket rate limiting (\#12), structured access logging (\#13), and OCSP stapling (\#14).

\section{Q-Queue: Universal Queue Fabric}

\subsection{Motivation}

Blockchain systems require low-latency message passing for consensus protocol messages, transaction mempool synchronization, and peer-to-peer gossip. Existing solutions present trade-offs:

\begin{table}[H]
\centering
\footnotesize
\begin{tabular}{llr}
\toprule
\textbf{System} & \textbf{Strength} & \textbf{Throughput} \\
\midrule
Chronicle Queue & Sub-$\mu$s latency & JVM only \\
Apache Kafka & Durability & 275 MB/s/node \\
Redpanda & Shared-nothing & 666 MB/s/node \\
VAST Broker & Disaggregated & 1.1 GB/s/node \\
\textbf{Q-Queue} & \textbf{All of above} & \textbf{$>$2 GB/s/node} \\
\bottomrule
\end{tabular}
\caption{Queue system comparison.}
\end{table}

Q-Queue unifies local IPC, persistent storage, and distributed messaging under a single API, leveraging Rust's zero-cost abstractions, io\_uring, and lock-free algorithms.

\subsection{Architecture}

\begin{figure}[H]
\centering
\begin{verbatim}
  +----------------------------------+
  |     Queue Client Library          |
  | (Producer/Consumer API)           |
  +------+----------+--------+-------+
         |          |        |
   +-----v---+ +---v----+ +-v--------+
   |  Local   | |Persist | |Distrib-  |
   |  IPC     | |ent     | |uted      |
   |(shmem)   | |(mmap)  | |(RDMA/TCP)|
   +----------+ +--------+ +----------+
\end{verbatim}
\caption{Q-Queue three-tier architecture.}
\end{figure}

\subsection{Phase 1: Lock-Free IPC Queue}

The core data structure is a bounded ring buffer with per-element versioning. Each slot contains an atomic version number that encodes both occupancy and generation:

\begin{lstlisting}
struct Slot<T> {
    version: AtomicU64,
    // odd = data ready, even = empty/writing
    data: UnsafeCell<MaybeUninit<T>>,
}

pub struct Queue<T> {
    slots: Box<[Slot<T>]>,
    producer_pos: AtomicUsize,
    consumer_pos: AtomicUsize,
}
\end{lstlisting}

\subsubsection{Enqueue Algorithm}

The producer atomically claims a slot using CAS on the version:

\begin{lstlisting}
fn enqueue(&self, value: T) -> Result<()> {
    loop {
        let pos = self.producer_pos
            .load(Ordering::Relaxed);
        let slot = &self.slots[
            pos % self.capacity];
        let v = slot.version
            .load(Ordering::Acquire);

        // Check: slot empty AND correct gen
        if v & 1 == 0
           && v / 2 == pos / self.capacity
        {
            if slot.version.compare_exchange(
                v, v + 1,
                Ordering::AcqRel,
                Ordering::Acquire
            ).is_ok() {
                // Write data
                unsafe {
                    (*slot.data.get())
                        .write(value);
                }
                // Publish: version becomes odd
                slot.version.store(
                    v + 2, Ordering::Release);
                self.producer_pos.fetch_add(
                    1, Ordering::Release);
                return Ok(());
            }
        }
        // Slot occupied or stale: retry
    }
}
\end{lstlisting}

\subsubsection{Dequeue Algorithm}

The consumer reads data when the version is odd (ready):

\begin{lstlisting}
fn dequeue(&self) -> Option<T> {
    let pos = self.consumer_pos
        .load(Ordering::Relaxed);
    let slot = &self.slots[
        pos % self.capacity];
    let v = slot.version
        .load(Ordering::Acquire);

    // Check: data ready AND correct gen
    if v & 1 == 1
       && v / 2 == pos / self.capacity
    {
        let value = unsafe {
            (*slot.data.get())
                .assume_init_read()
        };
        // Mark consumed: version becomes even
        slot.version.store(
            v + 1, Ordering::Release);
        self.consumer_pos.fetch_add(
            1, Ordering::Release);
        Some(value)
    } else {
        None // empty or not ready
    }
}
\end{lstlisting}

\subsubsection{Correctness Argument}

The version field serves dual purpose: (1)~the parity bit indicates slot state (even=empty, odd=ready), and (2)~the upper bits encode the generation, preventing ABA problems. A producer can only write to a slot when its version indicates empty \textit{and} the generation matches the expected cycle. This provides linearizable SPSC semantics without mutexes.

\subsubsection{Performance Analysis}

The critical path for enqueue is one CAS + one store (2 atomic operations). For dequeue, it is one load + one store (2 atomic operations). Cache-line contention is minimized by:

\begin{itemize}[nosep]
\item Separating producer and consumer positions onto different cache lines (64-byte padding)
\item Sequential slot access pattern exploits hardware prefetching
\item Avoiding false sharing via slot alignment
\end{itemize}

Target: $<$500~ns per message for same-machine transfers (vs. Chronicle Queue's $\sim$780~ns).

\subsection{Phase 2: Persistent Queue}

The persistent queue extends the lock-free design with memory-mapped files:

\begin{itemize}[nosep]
\item \textbf{Segment-based storage}: Pre-allocated file segments (default 256MB each) mapped via \texttt{mmap}
\item \textbf{Append-only writes}: CRC32 checksums per message for integrity
\item \textbf{Zero-copy reads}: Consumers map read-only views of completed segments
\item \textbf{Background compaction}: Expired segments reclaimed by a dedicated thread
\item \textbf{fsync policy}: Configurable per-message, per-batch, or per-second
\end{itemize}

Target: $>$10M messages/sec on NVMe storage.

\subsection{Phase 3: Distributed Queue}

The distributed queue uses a disaggregated architecture:

\begin{itemize}[nosep]
\item \textbf{Transport}: RDMA when available (InfiniBand, RoCE); fallback to TCP with io\_uring multishot receive
\item \textbf{Partitioning}: Consistent hashing with virtual nodes for balanced distribution
\item \textbf{Replication}: Erasure coding (Reed-Solomon) for space-efficient durability; configurable $k$-of-$n$ redundancy
\item \textbf{Protocol}: Optional Kafka wire protocol compatibility for ecosystem integration
\end{itemize}

\subsubsection{io\_uring Integration}

Q-Queue leverages io\_uring for both network and storage I/O:

\begin{lstlisting}
// Multishot receive: one SQE handles
// multiple incoming messages
let sqe = io_uring::opcode::
    RecvMsgMultishot::new(fd, msg, flags)
    .build()
    .user_data(conn_id);

// Cross-thread wakeup: zero syscalls
io_uring_prep_msg_ring(
    sqe, target_ring_fd,
    0, wakeup_data, 0
);
\end{lstlisting}

The io\_uring \texttt{msg\_ring} operation enables zero-syscall cross-thread signaling, replacing \texttt{eventfd} notifications.

Target: $>$2~GB/s per node, $>$5~GB/s for a 3-node cluster.

\subsection{Phase 4: SIMD Serialization}

For high-throughput scenarios, Q-Queue uses SIMD-accelerated serialization:

\begin{lstlisting}
#[cfg(target_arch = "x86_64")]
unsafe fn serialize_batch(
    buf: &mut [u8], values: &[i64]
) {
    // Process 4 i64s at once using AVX2
    let chunks = values.chunks_exact(4);
    for chunk in chunks {
        let v = _mm256_loadu_si256(
            chunk.as_ptr() as *const __m256i
        );
        _mm256_storeu_si256(
            buf.as_mut_ptr() as *mut __m256i,
            v
        );
    }
}
\end{lstlisting}

Runtime CPU feature detection ensures compatibility across hardware generations, with fallback to scalar paths.

\subsection{Integration with Q-NarwhalKnight}

Q-Queue integrates with the Q-NarwhalKnight consensus protocol at three levels:

\begin{enumerate}[nosep]
\item \textbf{Block propagation}: Gossipsub messages routed through local IPC queue between network layer and DAG-Knight consensus engine
\item \textbf{Transaction mempool}: Persistent queue ensures transaction durability across node restarts; zero-copy reads for block building
\item \textbf{P2P relay}: Distributed queue enables cross-node message forwarding without point-to-point connections (useful for Tor-routed nodes)
\end{enumerate}

\section{Production Lessons}

Deploying Q-NarwhalKnight on the Epsilon supernode with 400+ miners produced several insights relevant to both Q-Flux and Q-Queue:

\subsection{Connection Pooling is Non-Negotiable}

Setting \texttt{Connection: close} or \texttt{keepalive off} on a reverse proxy creates a new TCP connection per request. At 9,600~req/s, this consumed 42 of 48 CPU cores on connection setup/teardown alone, producing 87\% 503 error rates. Connection pooling with \texttt{max\_conns\_per\_host} caps reduced load from 121 to 24 and eliminated 503 errors entirely.

\subsection{Unbounded Spawns Cause OOM}

Any \texttt{tokio::spawn} that allocates significant memory must be gated by a semaphore. The block-pack handler crashed Epsilon 6+ times/hour because each syncing peer requested 500+ blocks ($\sim$50--150MB serialized), and multiple concurrent spawns allocated $>$10GB. Capping concurrent handlers to 4 limited worst-case allocation to 200MB.

\subsection{Atomic Ordering Selection}

For monotonic counters (metrics), \texttt{Relaxed} ordering is sufficient and avoids memory fence overhead. For visibility guarantees (shutdown flags), \texttt{SeqCst} ensures all cores see the flag immediately. For publish/subscribe patterns (lock-free queue), \texttt{Acquire/Release} pairs provide the minimum necessary ordering.

\section{Related Work}

\textbf{Reverse Proxies}: Envoy~\cite{envoy} uses an event-driven architecture with worker threads but relies on a shared L7 filter chain. HAProxy~\cite{haproxy} uses multi-process mode similar to Q-Flux but lacks native Rust integration. Pingora~\cite{pingora} pioneered Rust-based proxying but uses a single-process async model that limits TLS throughput.

\textbf{Queue Systems}: Disruptor~\cite{disruptor} introduced per-element sequencing for lock-free queues but is JVM-only. DPDK~\cite{dpdk} ring buffers achieve similar performance but require kernel bypass. Q-Queue combines Disruptor-style versioning with io\_uring for a pure-userspace solution.

\section{Conclusion}

Q-Flux and Q-Queue address the performance requirements of production blockchain infrastructure through shared-nothing architectures, lock-free algorithms, and zero-copy I/O paths. Q-Flux has been deployed serving 400+ concurrent miners on a 48-core supernode, replacing Caddy and eliminating the 87\% error rate. Q-Queue's lock-free ring buffer provides the foundation for sub-microsecond consensus message passing.

Both systems are open-source Rust crates within the Q-NarwhalKnight workspace. Issues \#1--\#14 in the project tracker describe ongoing and future work for community contribution via the collaborative development environment at \texttt{git://185.182.185.227/q-narwhalknight}.

\begin{thebibliography}{10}

\bibitem{qnk-yellow}
Q-NarwhalKnight Team, ``Q-NarwhalKnight Yellow Paper: DAG-BFT Consensus with Post-Quantum Cryptography,'' 2025.

\bibitem{seastar}
A. Rappoport et al., ``Seastar: High-performance server-side application framework,'' ScyllaDB, 2014.

\bibitem{block-pack-oom}
Q-NarwhalKnight Infrastructure, ``Block-Pack OOM Crash-Loop Fix (v9.1.9),'' Internal Report, March 2026.

\bibitem{envoy}
M. Klein, ``Envoy: C++ L7 proxy and communication bus,'' Lyft Engineering, 2016.

\bibitem{haproxy}
W. Tarreau, ``HAProxy: The Reliable, High Performance TCP/HTTP Load Balancer,'' 2001.

\bibitem{pingora}
Cloudflare, ``Pingora: Building HTTP proxy in Rust,'' Cloudflare Blog, 2024.

\bibitem{disruptor}
M. Thompson et al., ``Disruptor: High performance alternative to bounded queues,'' LMAX Exchange, 2011.

\bibitem{dpdk}
Intel, ``Data Plane Development Kit (DPDK),'' 2012.

\end{thebibliography}

\end{document}
